\section{Pressure Length and Elastic Moduli}
\label{sec:moduli}

\subsection{Pressure length}
From the Boerdijk--Coxeter rise $H=s/\sqrt{10}$ the critical linear strain is
\[
\varepsilon_c=\frac{1}{\sqrt{10}}.
\]
The \textbf{pressure length} quantifies residual-stress distortion of the cell.
Its magnitude is pure geometry:
\[
\ell=\varepsilon_c\,s=\frac{s}{\sqrt{10}}.
\]
Six combinatorial regimes relate residual pressure to $\ell$:

\begin{center}
\begin{tabular}{lll}
\toprule
Symbol & Condition & Meaning \\
\midrule
$\mathrm{pp}$ & both low & Highest compression \\
$\mathrm{P}\ell$ / $\mathrm{pl}$ & pressure dominates or strong compression & Strong compression \\
$\ell\mathrm{P}$ / $\mathrm{pL}$ & length dominates & Moderate compression / length-led \\
$\mathrm{PL}$ & pressure $=$ length (positive) & Balanced equality (positive) \\
$\mathrm{pl}$ & pressure $=$ length (negative) & Balanced equality (negative) \\
$\mathrm{PP}$ & both high & Highest stretch \\
\bottomrule
\end{tabular}
\end{center}

Primary comparison axis:
\[
\mathrm{P}\ell \;\longleftrightarrow\; \ell\mathrm{P}
\]

Full residual cycle:
\[
\mathrm{pp}\;\to\;\mathrm{pl}\;\to\;\mathrm{PL}\;\to\;\mathrm{pL}\;\to\;\mathrm{PP}\;\to\;\mathrm{Pl}\;\to\;\mathrm{pp}.
\]

\subsection{Elastic moduli from the tetrahedron}
\derived\ (ratios). Identifying residual stress with the load that produces the
geometric critical strains yields pure geometric ratios:
\begin{align*}
Y&=\sqrt{10}\,|\sigma_{\mathrm{res}}|,\\
K&=\frac{\sqrt{10}}{3}\,|\sigma_{\mathrm{res}}|,\\
G&=\frac{3}{\pi}\,|\sigma_{\mathrm{res}}|,
\end{align*}
where $Y$, $K$ and $G$ are Young-type, bulk and shear moduli respectively.
